% !TEX encoding = UTF-8 Unicode
\documentclass[a4paper,12pt]{article}
\usepackage[T2A]{fontenc}
\usepackage[utf8]{inputenc}
\usepackage[russian]{babel}
\usepackage{amsmath,amssymb,amsfonts}
\usepackage{bm}
\usepackage{graphicx}
\usepackage{hyperref}
\usepackage{cite}
\usepackage{geometry}
\geometry{left=2.5cm,right=2.5cm,top=2.5cm,bottom=2.5cm}
\usepackage{indentfirst}
\setlength{\parskip}{2mm}

\title{\textbf{AI-резонанс как техносингулярность:}\\
GRA-подход к аннулированию неустойчивостей в машинах, социуме и разуме}
\author{Олег Бит\\[2mm]
\small Московская городская телефонная сеть (МГТС), \texttt{\#OPEN\_TO\_WORK}\\[2mm]
\small \url{https://github.com/qqewq/GRA-Resonance-Nullification-Series}}
\date{22 июня 2026 г.}

\begin{document}
\maketitle

\begin{abstract}
Резонанс — универсальный механизм катастрофического усиления колебаний в физических, социальных и когнитивных системах. В настоящей работе представлена серия методов \textbf{генеративного аннулирования резонанса} (GRA) — открытый теоретико-экспериментальный фреймворк, позволяющий подавлять нежелательные резонансные моды за счёт целенаправленного введения противофазных управляющих сигналов или метаматериальных структур с отрицательными эффективными параметрами. Математический аппарат GRA обобщается на дискретные динамические системы, что открывает путь к контролю резонансных процессов в искусственном интеллекте, включая гипотетический сильный ИИ (AGI/ASI). Вводится концепция \textbf{AI-резонанса как одной из граней техносингулярности} и демонстрируется, как хиральные многоагентные архитектуры и управление в скрытом пространстве позволяют предотвратить неконтролируемое самовозбуждение. Полный программный код, данные экспериментов и теоретические выкладки доступны в репозитории \url{https://github.com/qqewq/GRA-Resonance-Nullification-Series}.
\end{abstract}

\section{Введение: от рушащихся мостов к коллапсирующим интеллектам}
Резонанс — это архитектура многих катастроф. От Токомакского моста до коллапсов финансовых рынков, от танкового сингулярного резонанса до потенциальных нестабильностей в системах искусственного интеллекта — механизм катастрофического усиления колебаний универсален. В настоящей работе мы предлагаем единый теоретический и инженерный фреймворк для подавления нежелательных резонансных мод.

\section{Математический аппарат GRA}
Рассмотрим линейную колебательную систему:
\begin{equation}
m\ddot{x} + c\dot{x} + kx = F(t)
\end{equation}
Резонансная частота $\omega_0 = \sqrt{k/m}$. Для подавления резонанса вводим управляющую силу $u(t)$, формируемую как:
\begin{equation}
u(t) = -K_p x(t) - K_d \dot{x}(t)
\end{equation}
что эквивалентно увеличению жёсткости и демпфирования системы.

\section{Резонатор с запаздыванием}
Альтернативный подход — использование запаздывающей обратной связи:
\begin{equation}
u(t) = \alpha x(t - \tau)
\end{equation}
При $\tau = \pi/\omega$ управляющий сигнал находится в противофазе с возбуждением, обеспечивая аннулирование резонанса.

\section{Метаматериальный подход}
Введение локальных резонаторов с массой $m_c$ и жёсткостью $k_c$ приводит к эффективной жёсткости:
\begin{equation}
K_{\text{eff}}(\omega) = k_h + \frac{k_c^2}{k_c - m_c \omega^2}
\end{equation}
При $\omega = \sqrt{k_c/m_c}$ знаменатель обращается в ноль, создавая антирезонанс — полное подавление колебаний.

\section{AI-резонанс и техносингулярность}
Распространим аналогию на дискретные динамические системы, моделирующие скрытые состояния нейронных сетей. Введём логический резонансный индекс (LRI):
\begin{equation}
\text{LRI}_t = 1 - \cos\theta_t, \quad \cos\theta_t = \frac{\mathbf{h}_t \cdot \mathbf{h}_{t+1}}{\|\mathbf{h}_t\| \|\mathbf{h}_{t+1}\|}
\end{equation}
Высокое значение LRI указывает на резонансное самовозбуждение в скрытом пространстве.

\section{Хиральные многоагентные архитектуры}
Для предотвращения неконтролируемого роста в многоагентных системах предлагается использовать хиральные пары агентов с противоположной полярностью, взаимно гасящие резонансные моды.

\section{Заключение}
GRA-фреймворк предоставляет единый математический и программный инструментарий для подавления резонансных неустойчивостей в широком спектре систем — от механических конструкций до искусственного интеллекта.

\end{document}
